\documentclass[a4paper]{extarticle}
%\documentclass[12pt, a4paper]{article}
% Кодировки и язык
\usepackage[T2A]{fontenc}          % Кириллическая кодировка шрифтов
\usepackage[utf8]{inputenc}        % Кодировка исходного файла — UTF-8
\usepackage[russian]{babel}        % Поддержка русского языка

% Основные пакеты
\usepackage{graphicx}              % Работа с изображениями
\usepackage{amsmath}               % Математические формулы
\usepackage{caption}
\usepackage{indentfirst}           % Отступ в первом абзаце раздела
\usepackage{fancyvrb}              % Расширенный режим Verbatim
\usepackage{fvextra}               % Поддержка breaklines и UTF-8 в Verbatim
\usepackage{comment}               % Блоки комментариев
\usepackage{subfiles}              % Модульная сборка документа
\usepackage{enumitem}              % Гибкая настройка списков
\usepackage[skip=0pt]{parskip}     % Убирает лишние отступы между абзацами
\usepackage{float}
\usepackage{makecell}
\usepackage{hyperref}
% Гиперссылки — всегда подключать ПОСЛЕДНИМ из основных пакетов

\usepackage{geometry}
\usepackage{afterpage}
\usepackage{scrlayer-scrpage}
\usepackage{pdflscape}
\usepackage{pdfpages}

% === Настройка листингов (без плавающих объектов) ===
% Используем ручной счётчик — это надёжно для длинных кодов с русским текстом
\newcounter{listing}
\renewcommand{\thelisting}{\arabic{listing}}
\newcommand{\listingcaption}[1]{%
  \par\refstepcounter{listing}%
  \noindent\textbf{Листинг~\thelisting.} #1\par\nobreak\vspace{0.4em}%
}

% === Путь к изображениям ===
\graphicspath{{./images/}}

% === Геометрия страницы ===
\textheight=24cm    % Высота текстового блока
\textwidth=16cm     % Ширина текстового блока
\oddsidemargin=0pt  % Отступ от левого края (для односторонней печати)
\topmargin=-1.5cm   % Верхний отступ
\parindent=24pt     % Абзацный отступ
\flushbottom        % Выравнивание текста по нижнему краю страницы